\subsolution{Переменный (?) ток}{3.5}

\textbf{Внимание:} стандартным шрифтом описан исходный авторский метод; {\color{violet}фиолетовым шрифтом описан альтернативный метод.}

\criterion[type=remark]
{Критерии за альтернативный метод эквивалентны критериям за описание фазовых диаграмм \textbf{(ФД)}. Общий балл ставится по наибольшему между двумя методами внутри каждого из двух частей задачи. Переноса ошибки на численные расчёты нет.}

Цепь выглядит так, словно на $LC$-участок действует периодический импульс $U(t)$, показанный на Рис. \figref{1}. Конечно, эта схема учитывает идеализированную подачу внешнего напряжения, при которой $U(t)$ не ``забирает'' ток к себе --- иначе конденсатор мгновенно бы имел заряд $q(t)=CU(t)$!

\cpic{0.42}{figures/1b.2.png}[Напряжение в виде коротких импульсов длительностью $\tau$ и периодом $T$]

Период колебаний $LC$-контура равен известной из школьной программы формулы Томпсона
\begin{eq}[points=0.2, tail={; принимается $\omega=\sqrt{LC}$}]
    T_0 = 2\pi\sqrt{LC}.
\end{eq}

Обратим внимание, что численно

\begin{eq}[points=0.1]
    \frac T{T_0} = \frac14
\end{eq}

с очень высокой точностью. Это пригодится нам для решения задачи дальше. Действие ``импульса'' $U_0\tau$ выражается в резком скачке тока $I$ на $LC$-участке. Действительно, мы имеем по второму правилу Кирхгофа выражение

\begin{eq}[points=0.1, tail={; принимается $U(t)\Leftrightarrow U_0\Leftrightarrow0$ в правой части}]
    L\cdot\frac{dI}{dt} + \frac qC = U(t),
\end{eq}

и, проинтегрировав за промежуток $\tau$, мы видим, что ввиду конечности силы тока мы \textbf{не} можем иметь резкий скачок заряда $q$, так что интегрирование \eqref{1} упрощается до

\begin{eq}[points=0.2]
    \Delta I = \frac{U_0\tau}L.
\end{eq}

Самый удобный в данный задаче метод заключается в рассмотрении эволюции данной системы в фазовом пространстве $(q, I)$. Поскольку колебания напряжения на конденсаторе и индуктивности являются чисто синусоидальными (за исключением резких скачков) с относительным сдвигом фаз $\pi/2$, ``точка состояния'' системы $(q, I)$ будет двигаться вдоль окружности по часовой стрелке. Это видно из выражения для энергии системы
$$
E=\frac{LI^2}{2} + \frac{q^2}{2C},
$$
что будет уравнением окружности, если ось тока привести в единицы $I/\sqrt{LC}$. За промежуток $T$, поскольку это четверть периода колебаний $T_0$, фазовое состояние повернётся на 90 градусов по часовой стрелке; $\Delta I$ будет иметь смысл вертикального вектора, позволяющего фазовой точке ``перескочить'' с одной окружности на другую. Графическое решение показано на Рис. \figref{2}. Фазовое состояние при $t=0$ находилось на $(0,0)$; затем точка шла по красной $\to$ синей $\to$ оранжевой $\to$ зелёной траекториям и вернулась обратно к состоянию $(0,0)$.
\criterion{\textbf{Метод фазовых диаграмм:} (следующие 5 пунктов эквивалентны \eqref{6}--\eqref{7})}
\criterion[points=0.1]{\textbf{(ФД)} Движение фазового состояния по окружности \textbf{и} по часовой стрелке}
\criterion[points=0.1]{\textbf{(ФД)} Нормализация $I/\sqrt{LC}$ или эквивалент}
\criterion[points=0.1]{\textbf{(ФД)} Каждый импульс --- это вертикальный вектор перемещения точки}
\criterion[points=0.1]{\textbf{(ФД)} Ввиду \eqref{2} происходит поворот на $\pi/2$ между каждым импульсом}
\criterion[points=0.3]{\textbf{(ФД)} Верно нарисовано движение точки (как на Рис. \figref{2})}
Оттуда сразу заключаем, что энергия в момент $t=3T+\tau$ равна

\begin{eq}[points=0.4, prefix={Ответ к \textbf{(a)}}, suffix={; без обоснования не принимается}]
E=0.
\end{eq}

{\color{violet} \textit{Альтернативно,} рассмотрим ``в лоб'' эволюцию заряда и тока от времени. После импульса напряжение источника отсутствует, и система совершает свободные колебания с частотой $\omega$, получаемого из \eqref{1}:

\criterion[type=remark]{\textbf{Альтернативный метод:} (вместо \textbf{ФД})}

$$
\omega=\frac1{\sqrt{LC}}.
$$

В момент $t=\tau$ начальные условия $q(\tau)=0$, $I(\tau)=\Delta I$, так что для $\tau <t<T$,

\begin{eq}[points=0.2, type=alternative]
    q(t)=\frac{\Delta I}{\omega}\sin(\omega t).
\end{eq}

Через время $T$, поскольку это четверть периода колебаний, сила тока обнуляется, а заряд становится максимальным и равным $\Delta I/\omega$. После второго импульса сила тока становится равной $I(T+\tau) = \Delta I$, так что уравнением дальнейших колебаний для $q(t)$ при $T+\tau<t<2T$ будет

\begin{eq}[points=0.5, type=alternative]
    q(t)=\frac{\Delta I}\omega\ab(\sin\omega(t-T)+\cos\omega(t-T))=\frac{\sqrt2\Delta I}\omega\sin\ab(\omega (t-T)+\frac\pi4).
\end{eq}

К моменту $t=2T$ мы видим, что $q(2T)=q(T+\tau)=\Delta I/\omega$, а $I(2T)=-I(T+\tau)=-\Delta I$, так что после третьего импульса $I(2T+\tau)=0$. Система перешла в состояние отсутствия тока и максимума заряда для процесса $2T+\tau<t<3T$ --- следовательно, спустя четверть периода она перейдёт в состояние максимума модуля тока $I=-\Delta I$ и отсутствия заряда. После четвёртого импульса мы можем заключить о \eqref{5}.
}
Максимальный заряд равен точке пересечения синей траектории с горизонтальной осью на внешней окружности. Радиус внешней окружности, как видно, равен $\sqrt2$ от длины вертикального вектора $\Delta I\sqrt{LC}$. Это сразу нас приводит к ответу

\begin{eq}[points=0.3, tail={; принимается факт, что внешняя окружность в $\sqrt2$ раза больше внутренней, либо факт, что энергия после второго импульса в $2^{1/4}$ раз больше}]
    q = \sqrt2\cdot\Delta I\sqrt{LC}.
\end{eq}

\criterion[points=-0.2, type=penalty]{Забыт множитель $\sqrt{LC}$}

{\color{violet} Если вернёмся к выражению \eqref{7}, то можно заключить о том же самом при подстановке $t=1.5T$.} Имея выражение \eqref{4}, получаем окончательно

\begin{eq}[points=0.3, numpoints=0.1, prefix={Ответ к \textbf{(a)}}]
    q = U_0\tau\sqrt\frac{2C}{L}=2.09\pu{мКл}.
\end{eq}

Решение пункта \textbf{(b)} методом фазовых диаграмм тоже достаточно прямолинейно и быстро. Мы видим, что $q_*$ диктует радиус окружности в фазовом пространстве. Поскольку энергия системы не меняется из-за подачи импульса, вертикальный вектор $\Delta I$ должен являться хордой данной окружности с секторным углом $\pi/2$. Поскольку в момент $t=0$ точка находилась на $(q_*,0)$, ей нужно пройти дугу углом $\pi/4$ чтобы прийти к моменту скачка. {\color{violet} Это так же видно ``просмотром назад во времени'': в момент $t_*-T$ состояние системы должно повторять состояние системы в $t_*+\tau$.} А значит,

\criterion{\textbf{Метод фазовых диаграмм:} (следующий 1 пункт эквивалентен \eqref{11})}
\criterion[points=0.3]{\textbf{(ФД)} Указано, что $I/\sqrt{LC}$ должен опираться хордой под углом 90 градусов на окружность $q_*$}

\begin{eq}[points=0.3, prefix={Ответ к \textbf{(b)}}, tail={; за отсутствие численного значения штрафа нет}]
    t_* = \frac T2=125\pu{мс}.
\end{eq}

Теперь к $q_*$. {\color{violet} Можно заметить, что работа источника $A = U_0\Delta q$ равняется нулю ввиду сохранения полной энергии $E$. Это возможно только в том случае, если сохраняется модуль силы тока $I$ и меняется его знак: начиная с выражения $q(t)=q_*\cos\omega t$, имеем}

\criterion[type=remark]{\textbf{Альтернативный метод:} вместо \textbf{ФД}}

\begin{eq}[points=0.3, type=alternative]
-\omega q_*\sin\omega t_* = -\frac{\Delta I}{2},
\end{eq}

{\color{violet} поскольку после скачка сила тока станет $\Delta I/2$. Подставив \eqref{10} в \eqref{11}, можно решить относительно $q_*$,} или же, нарисовав как на Рис. \figref{2}, геометрией заключаем

\begin{eq}
    [points=0.4, numpoints=0.1, prefix={Ответ к \textbf{(b)}}, tail={; балл не ставится если забыт множитель $\sqrt{LC}$}]
    q_* = \frac1{\sqrt2}\cdot\Delta I\sqrt{LC}=U_0\tau\sqrt\frac C{2L}=1.04\pu{мКл}.
\end{eq}

\cpic{1}{figures/1b.3}[Фазовые диаграммы для пунктов \textbf{(a)} (слева) и \textbf{(b)} (справа)]
